\subsection{可数型可度量化空间}

定义 4. 可度量化空间称为可数型的（或可分的），如果它的拓扑具有可数基。

显然，可数型可度量化空间的每个子空间仍是可数型的。乘积空间拓扑基的定义（第一章，§ 4，第 1 号）与第 4 号定理 1 的推论 2 表明，可数个可数型可度量化空间的乘积是可数型可度量化空间。同样，可数个可数型可度量化空间的和是可数型可度量化空间（第 5 号）。

命题 12. 若 X 是可度量化拓扑空间，则下列陈述等价：

\begin{enumerate}
\def\labelenumi{\alph{enumi})}
\item
  X 是可数型的。
\item
  X 有可数稠密子集。
\item
  X 同胚于方体 IN 的一个子空间，其中 I 是 R 中的区间 {[}o, i{]} 。
\end{enumerate}

由前面的注记，显然 c) 蕴含 a)；a) 蕴含 b)，因为若 \((\mathbf{U}_{n})\) 是 X 的拓扑的一个可数基，且 \(a_{n}\) 是 \(U_{n}\) 的一个点，则诸 \(a_{n}\) 构成 X 的一个稠密子集。最后证明 b) 蕴含 c)。设 \((a_{n})\) 是 X 的点的稠密序列，且对每个 \(x \in X\) ，令 \(\varphi(x)\) 为点 \((d(x, a_{n}))_{n \in \mathbb{N}}\) —— \(I^{N}\) 的点 ——（d 是与 X 的拓扑相容、且关于它 X 的直径 \(\leqslant 1\) 的一个度量）。我们将证明 \(\varphi\) 是 X 到 \(I^{N}\) 的一个子空间上的同胚。事实上，\(\varphi\) 是连续的，因为每个函数 \(x \to d(x, a_{n})\) 都连续；且 \(\varphi\) 是单射，因为 X 的每个点都是 \((a_{n})\) 的某个子序列的极限（命题 8）。设 B 是 X 中以 \(x_{0}\) 为中心、以 r 为半径的一个球，并设 n 是使

\[
\left| d (x _ {0}, a _ {n}) <   \frac {1}{3} r. \right|
\]

\(\varphi\) 把集 \(W\) —— 由满足以下条件的点 \(x\in X\) 组成的集 ——

\[
\left| d (x _ {0}, a _ {n}) - d (x, a _ {n}) \right| <   \frac {1}{3} r
\]

映为按定义是 \(\varphi(x_0)\) 在 \(\varphi(\mathbf{X})\) 中的一个邻域。但对每个 \(x \in W\) 有 \(d(x, a_n) < d(x_0, a_n) + \frac{1}{3} r < \frac{2}{3} r\) ，由此得

\[
d (x, x _ {0}) \leqslant d (x _ {0}, a _ {n}) + d (x, a _ {n}) <   r,
\]

这表明 W 是包含在 B 中的 \(x_{0}\) 的一个邻域，因而 \(\varphi\) 是 X 到 \(\varphi(\mathbf{X})\) 上的一个同胚。

\pandocbounded{\includegraphics[keepaspectratio]{images/part001_a8399183.jpg}}

注意，对任意拓扑空间 X 而言，性质 b) 不一定蕴含可数基的存在性，即使 X 是紧的且 X 的每个点都有可数的基本邻域系（习题 13；参看第一章，§ 1，习题 7）。

命题 13. 设 X 是具有可数基 \((\mathbf{U}_n)\) 的拓扑空间；则对 X 的每个开覆盖 \((\mathbf{V}_t)_{i \in \mathbf{I}}\) ，存在 I 的可数子集 J，使得 \((\mathbf{V}_t)_{i \in \mathbf{J}}\) 是 X 的一个覆盖。

设 H 是由使 \(U_{n}\) 至少包含在某个 \(V_{t}\) 中的指标 n 组成的 N 的子集；序列 \((\mathrm{U}_{n})_{n\in\mathbf{H}}\) 是 X 的一个覆盖，因为每个点 \(x\in X\) 属于某个 \(V_{t}\) ，且由于 \(V_{t}\) 是开的，存在指标 n 使得 \(x\in U_{n}\subset V_{t}\) 。因而存在 H 到 I 中的一个映射 \(\psi\) ，使得 \(U_{n}\subset V_{\psi(n)}\) 对每个 \(n\in H\) 成立；取 \(J=\psi(H)\) ，它是可数的，命题得证。

